\documentclass[../../dis.tex]{subfiles}
\begin{document}
\chapter{L5 - Access Methods: Hashing and B\texorpdfstring{$^+$}{+}-Trees}\label{chap:l5}

\section{Where access methods fit}
The DBMS execution engine reads and writes \textbf{pages}. To avoid scanning every page for every query, the DBMS uses \textbf{access methods}: physical data structures that help it find the right page or record quickly.

At this layer, two large families dominate:
\begin{dashlist}
    \item \textbf{tree-based structures}: ordered access paths,
    \item \textbf{hash-based structures}: unordered access paths.
\end{dashlist}

This difference is the key idea of the lecture:
\begin{dashlist}
    \item if the workload needs mainly \textbf{equality search}, hashing is often attractive,
    \item if the workload needs \textbf{equality and range search}, ordered trees are the standard choice.
\end{dashlist}

\section{How an index points to table data}
\begin{concept}{Data entry \(k^*\)}
The information stored in the index for search-key value \(k\). It tells the DBMS how to reach the matching table record or records.
\end{concept}

There are three classic representations for a data entry:

\begin{dbtable}{Three alternatives for index data entries}{C{1.3cm}L{3.0cm}L{5.2cm}}
\textbf{Alt.} & \textbf{Form of entry} & \textbf{Meaning} \\ \hline
1 & actual data record with key \(k\) & The index entry is the record itself. The index and the data file are physically the same structure. \\ \hline
2 & \(\langle k, \text{RID} \rangle\) & The index stores the key and one pointer to one matching record. \\ \hline
3 & \(\langle k, \text{list of RIDs} \rangle\) & The index stores one key and a list of pointers to all matching records. \\
\end{dbtable}

The important point is that this choice is \textbf{orthogonal} to the indexing technique:
\begin{dashlist}
    \item a \textbf{hash index} can use Alternative 1, 2, or 3,
    \item a \textbf{tree index} can also use Alternative 1, 2, or 3.
\end{dashlist}

\subsection{Reading the three alternatives}
\begin{dashlist}
    \item \textbf{Alternative 1}: the leaf level stores the full records. Very fast because no extra jump from index entry to data record is needed.
    \item \textbf{Alternative 2}: one pair \(\langle k, \text{RID} \rangle\) per matching record. Simple, works well even when the search key is not unique.
    \item \textbf{Alternative 3}: one key with a full RID list. Saves space when many records share the same key value.
\end{dashlist}

\subsection{Why Alternatives 2 and 3 are often preferred}
Suppose an employee file is itself hashed on \texttt{age}. Then the buckets store the \textbf{actual employee records}: this is Alternative 1.

Now suppose we also want an index on \texttt{salary}. We usually do \emph{not} want to duplicate the full records inside a second structure. Instead, we create another access path that stores either:
\begin{dashlist}
    \item \(\langle \texttt{salary}, \text{RID} \rangle\) pairs, or
    \item \(\langle \texttt{salary}, \text{list of RIDs} \rangle\) entries.
\end{dashlist}

This leads to the standard trade-offs:
\begin{dashlist}
    \item Alternatives 2 and 3 are easier to maintain than Alternative 1.
    \item If a file has several indexes, only \textbf{one} of them can use Alternative 1, because the data records can only be physically organized one way at a time.
    \item Alternative 3 is usually more compact than Alternative 2 when many duplicates exist.
    \item The downside of Alternative 3 is that entries become \textbf{variable-sized}, even when the search key has fixed length. If one key has a very long RID list, that entry may span several pages.
\end{dashlist}

\section{Hash indexes versus tree indexes}
Hashing and trees solve different problems.

\begin{concept}{Hash-based index}
Uses a hash function \(h(k)\) to send a search-key value to one bucket. Because the buckets are not ordered by key value, hashing is best for equality search.
\end{concept}

\begin{concept}{Tree-based index}
Keeps entries in sorted search-key order. In practice, the standard disk-oriented tree structure is the \mbox{$B^+$-tree}.
\end{concept}

\begin{dbtable}{Hash indexes and tree indexes}{L{2.7cm}L{3.9cm}L{3.6cm}}
\textbf{Structure} & \textbf{Best at} & \textbf{Main limitation} \\ \hline
Hash index & equality predicates such as \texttt{age = 44} & no natural order, so range predicates are poor \\ \hline
\mbox{$B^+$-tree} index & equality, range, ordered scans, sequential access & usually a little more overhead for pure equality lookups \\
\end{dbtable}

\subsection{Why range search needs order}
Consider the query:
\begin{sql}
SELECT *
FROM Student
WHERE gpa > 3.0
\end{sql}

If the data file were physically sorted on \texttt{gpa}, the DBMS could:
\begin{dashlist}
    \item do a binary search to find the first matching student,
    \item then scan forward to collect the rest.
\end{dashlist}

The problem is that keeping the \emph{whole data file} sorted is expensive: insertions and deletions constantly disturb the order.

The simple idea is therefore:
\begin{dashlist}
    \item keep a smaller \textbf{index file} in sorted order,
    \item search the index first,
    \item then use it to reach the data pages.
\end{dashlist}

This is the basic intuition behind the \mbox{$B^+$-tree}: keep the \emph{index} ordered, not the whole relation.

\section{B\texorpdfstring{$^+$}{+}-trees}
\begin{concept}{B\texorpdfstring{$^+$}{+}-tree}
A self-balancing, ordered tree structure optimized for page-based storage. Search, insertion, deletion, and sequential access all work efficiently, typically in \(O(\log_F N)\), where \(F\) is the fanout and \(N\) is the number of leaf nodes.
\end{concept}

\begin{concept}{Fanout}
The maximum number of child pointers an internal node can hold. Large fanout is what makes disk-based trees shallow.
\end{concept}

\subsection{Why B\texorpdfstring{$^+$}{+}-trees are so widely used}
\mbox{$B^+$-trees} and their variants are the preferred tree structures for external storage because they match the way disks and SSDs work:
\begin{dashlist}
    \item one node is typically one \textbf{page},
    \item each internal node stores many separator keys and child pointers,
    \item high fanout means very small tree height,
    \item all leaves are linked in sorted order, which makes range scans efficient.
\end{dashlist}

The broader family contains several variants, such as \textbf{B-tree}, \mbox{$B^+$-tree}, \textbf{B*-tree}, \textbf{B-link tree}, and \textbf{B\(\varepsilon\)-tree}. In this lecture, the main focus is the \mbox{$B^+$-tree}.

\subsection{Core properties}
A \mbox{$B^+$-tree} is an \(n\)-way search tree with the following properties:
\begin{dashlist}
    \item nodes come in three roles: \textbf{root}, \textbf{inner}, and \textbf{leaf},
    \item the tree is \textbf{perfectly balanced}: every leaf is at the same depth,
    \item every node except the root is at least \textbf{half full},
    \item an inner node has between \(\lceil n/2 \rceil\) and \(n\) children,
    \item if an inner node has \(k\) children, it has \(k-1\) separator keys,
    \item a leaf node has between \(\left\lceil (n-1)/2 \right\rceil\) and \(n-1\) keys,
    \item the root is the one exception: it may contain fewer entries than ordinary internal nodes.
\end{dashlist}

These conditions are exactly what keep the tree balanced and keep the height small.

\begin{schema}
{\begin{tikzpicture}[
    x=1cm,y=0.9cm,
    every node/.style={font=\small},
    >=Latex,
    inode/.style={draw=black!60, rounded corners=3pt, line width=0.85pt, fill=blue!10, minimum height=0.82cm, inner xsep=12pt},
    leaf/.style={draw=black!60, rounded corners=3pt, line width=0.85pt, fill=blue!7, minimum height=0.82cm, inner xsep=10pt},
    link/.style={->, thick, draw=black!55},
    leaflink/.style={<->, thick, draw=red!65!black}
]
    \node[inode] (root) at (6.0,5.2) {20};
    \node[inode] (lefti) at (3.1,3.3) {10};
    \node[inode] (righti) at (8.9,3.3) {35};

    \node[leaf] (l1) at (0.9,1.2) {6};
    \node[leaf] (l2) at (3.5,1.2) {10};
    \node[leaf] (l3) at (6.6,1.2) {20 \quad 31};
    \node[leaf] (l4) at (10.0,1.2) {38 \quad 44};

    \draw[link] (root.south west) -- node[pos=0.55, above left, font=\scriptsize] {$< 20$} (lefti.north);
    \draw[link] (root.south east) -- node[pos=0.55, above right, font=\scriptsize] {$\geq 20$} (righti.north);

    \draw[link] (lefti.south west) -- node[pos=0.55, above left, font=\scriptsize] {$< 10$} (l1.north);
    \draw[link] (lefti.south east) -- node[pos=0.55, above right, font=\scriptsize] {$\geq 10$} (l2.north);

    \draw[link] (righti.south west) -- node[pos=0.55, above left, font=\scriptsize] {$< 35$} (l3.north);
    \draw[link] (righti.south east) -- node[pos=0.55, above right, font=\scriptsize] {$\geq 35$} (l4.north);

    \draw[leaflink] (l1.east) -- (l2.west);
    \draw[leaflink] (l2.east) -- (l3.west);
    \draw[leaflink] (l3.east) -- (l4.west);

    \node[anchor=west, font=\scriptsize] at (11.8,5.2) {root node};
    \node[anchor=west, font=\scriptsize] at (11.8,3.3) {inner nodes};
    \node[anchor=west, font=\scriptsize] at (11.8,1.2) {leaf nodes};
    \node[font=\scriptsize] at (6.0,0.15) {leaf links preserve key order for range scans};
\end{tikzpicture}
}\par}
\end{schema}

In a \mbox{$B^+$-tree}:
\begin{dashlist}
    \item \textbf{internal nodes} store only separator keys and child pointers,
    \item \textbf{leaf nodes} store the actual data entries,
    \item \textbf{sibling pointers} connect leaves from left to right, so after the first match the DBMS can keep scanning sequentially.
\end{dashlist}

The data entry stored in each leaf can still follow Alternative 1, 2, or 3. The tree determines \emph{how to find} the entry; the alternative determines \emph{what the entry contains}.

\subsection{How search works}
Search always starts at the root. Key comparisons choose one child pointer at each internal node until the search reaches a leaf.

\begin{schema}
{\begin{tikzpicture}[
    x=1cm,y=0.9cm,
    every node/.style={font=\small},
    >=Latex,
    inode/.style={draw=black!60, rounded corners=3pt, line width=0.85pt, fill=blue!10, minimum height=0.84cm, inner xsep=10pt},
    leaf/.style={draw=black!60, rounded corners=3pt, line width=0.85pt, fill=blue!7, minimum height=0.84cm, inner xsep=8pt},
    link/.style={->, thick, draw=black!55},
    leaflink/.style={<->, thick, draw=red!65!black}
]
    \node[inode] (root) at (7.0,4.8) {13 \quad 17 \quad 24 \quad 30};

    \node[leaf] (l1) at (0.8,1.5) {2 \quad 3 \quad 5 \quad 7};
    \node[leaf] (l2) at (4.0,1.5) {14 \quad 16};
    \node[leaf] (l3) at (7.1,1.5) {19 \quad 20 \quad 22};
    \node[leaf] (l4) at (10.4,1.5) {24 \quad 27 \quad 29};
    \node[leaf] (l5) at (14.1,1.5) {33 \quad 34 \quad 38 \quad 39};

    \draw[link] (root.south west) -- (l1.north);
    \draw[link] ($(root.south west)!0.38!(root.south east)$) -- (l2.north);
    \draw[link] (root.south) -- (l3.north);
    \draw[link] ($(root.south west)!0.74!(root.south east)$) -- (l4.north);
    \draw[link] (root.south east) -- (l5.north);

    \draw[leaflink] (l1.east) -- (l2.west);
    \draw[leaflink] (l2.east) -- (l3.west);
    \draw[leaflink] (l3.east) -- (l4.west);
    \draw[leaflink] (l4.east) -- (l5.west);

    \node[font=\scriptsize] at (7.0,0.25) {exact search stops in one leaf; range search continues through sibling pointers};
\end{tikzpicture}
}\par}
\end{schema}

The example above shows three typical cases:

\begin{dbtable}{Reading searches in a B\texorpdfstring{$^+$}{+}-tree}{L{2.2cm}L{3.8cm}L{4.3cm}}
\textbf{Query} & \textbf{Leaf reached} & \textbf{What we learn} \\ \hline
\texttt{find 5} & \([2,3,5,7]\) & The key is present in that leaf. \\ \hline
\texttt{find 15} & \([14,16]\) & The key is not in the tree, because it would have to appear in that leaf between 14 and 16. \\ \hline
find keys \(\geq 24\) & \([24,27,29]\) & Start at the first matching leaf, then follow sibling pointers to the right. \\
\end{dbtable}

This is why \mbox{$B^+$-trees} support both:
\begin{dashlist}
    \item \textbf{equality search}: descend to one leaf and check it,
    \item \textbf{range search}: descend once, then scan leaf pages in order.
\end{dashlist}

\subsection{Insertion example: insert 8}
\begin{concept}{Leaf split}
If the target leaf is full, the DBMS allocates a new leaf and redistributes the entries. Then it inserts a separator key into the parent. If the parent also overflows, the split propagates upward.
\end{concept}

Suppose we insert key \(8\) into the tree shown below.

\begin{schema}
{\begin{tikzpicture}[
    x=0.92cm,y=0.84cm,
    every node/.style={font=\small},
    >=Latex,
    inode/.style={draw=black!60, rounded corners=3pt, line width=0.85pt, fill=blue!10, minimum height=0.82cm, inner xsep=9pt},
    leaf/.style={draw=black!60, rounded corners=3pt, line width=0.85pt, fill=blue!7, minimum height=0.82cm, inner xsep=7pt},
    fullleaf/.style={draw=orange!80!black, rounded corners=3pt, line width=1.05pt, fill=orange!12, minimum height=0.82cm, inner xsep=7pt},
    newleaf/.style={draw=green!45!black, rounded corners=3pt, line width=1.0pt, fill=green!12, minimum height=0.82cm, inner xsep=7pt},
    link/.style={->, thick, draw=black!55},
    leaflink/.style={<->, thick, draw=red!65!black},
    promote/.style={->, dashed, very thick, draw=red!70!black}
]
    \node[font=\bfseries] at (6.8,7.2) {Before};
    \node[inode] (r0) at (6.8,6.0) {13 \quad 17 \quad 24};

    \node[fullleaf] (a0) at (0.8,3.7) {2 \quad 3 \quad 5 \quad 7};
    \node[leaf] (b0) at (4.2,3.7) {14 \quad 16};
    \node[leaf] (c0) at (8.0,3.7) {19 \quad 20 \quad 22 \quad 23};
    \node[leaf] (d0) at (11.9,3.7) {24 \quad 27 \quad 29};

    \draw[link] (r0.south west) -- (a0.north);
    \draw[link] ($(r0.south west)!0.34!(r0.south east)$) -- (b0.north);
    \draw[link] ($(r0.south west)!0.66!(r0.south east)$) -- (c0.north);
    \draw[link] (r0.south east) -- (d0.north);

    \draw[leaflink] (a0.east) -- (b0.west);
    \draw[leaflink] (b0.east) -- (c0.west);
    \draw[leaflink] (c0.east) -- (d0.west);

    \node[font=\scriptsize, text=orange!80!black] at (0.8,2.9) {target leaf is full};

    \node[font=\bfseries] at (6.8,1.8) {After};
    \node[inode] (r1) at (6.8,0.6) {5 \quad 13 \quad 17 \quad 24};

    \node[leaf] (a1) at (-0.4,-1.8) {2 \quad 3};
    \node[newleaf] (b1) at (2.5,-1.8) {5 \quad 7 \quad 8};
    \node[leaf] (c1) at (5.9,-1.8) {14 \quad 16};
    \node[leaf] (d1) at (9.8,-1.8) {19 \quad 20 \quad 22 \quad 23};
    \node[leaf] (e1) at (13.7,-1.8) {24 \quad 27 \quad 29};

    \draw[link] (r1.south west) -- (a1.north);
    \draw[link] ($(r1.south west)!0.24!(r1.south east)$) -- (b1.north);
    \draw[link] ($(r1.south west)!0.48!(r1.south east)$) -- (c1.north);
    \draw[link] ($(r1.south west)!0.74!(r1.south east)$) -- (d1.north);
    \draw[link] (r1.south east) -- (e1.north);

    \draw[leaflink] (a1.east) -- (b1.west);
    \draw[leaflink] (b1.east) -- (c1.west);
    \draw[leaflink] (c1.east) -- (d1.west);
    \draw[leaflink] (d1.east) -- (e1.west);

    \draw[promote] (b1.north) to[bend left=18] node[midway, right, font=\scriptsize] {copy up 5} (r1.south);

    \node[font=\scriptsize, text=green!40!black] at (2.5,-2.65) {new right leaf after split};
\end{tikzpicture}
}\par}
\end{schema}

The steps are:
\begin{dashlist}
    \item search from the root and reach the first leaf, because \(8 < 13\),
    \item the target leaf already contains \([2,3,5,7]\), so it is full,
    \item allocate one new leaf and redistribute the values evenly,
    \item after inserting \(8\), the two leaves become \([2,3]\) and \([5,7,8]\),
    \item copy separator key \(5\) into the parent, because \(5\) is now the first key of the new right leaf,
    \item update sibling pointers so leaf-level scans still work left to right.
\end{dashlist}

In this example, the parent still has room, so the insertion stops there. If the parent were also full, the same split idea would continue upward, possibly all the way to the root.

\subsection{Insertion can propagate to the root}
\begin{concept}{Root split}
If an insertion reaches an overflowing root, the DBMS splits the root and creates a new root above it. This is the only case where the height of a \mbox{$B^+$-tree} increases.
\end{concept}

Now insert key \(21\) into the current tree. The target leaf is \([19,20,22,23]\), which is already full.

\begin{schema}
{\resizebox{0.98\textwidth}{!}{\begin{tikzpicture}[
    x=0.88cm,y=0.80cm,
    every node/.style={font=\small},
    >=Latex,
    inode/.style={draw=black!60, rounded corners=3pt, line width=0.85pt, fill=blue!10, minimum height=0.82cm, inner xsep=8pt},
    leaf/.style={draw=black!60, rounded corners=3pt, line width=0.85pt, fill=blue!7, minimum height=0.82cm, inner xsep=7pt},
    fullleaf/.style={draw=orange!80!black, rounded corners=3pt, line width=1.05pt, fill=orange!12, minimum height=0.82cm, inner xsep=7pt},
    newleaf/.style={draw=green!45!black, rounded corners=3pt, line width=1.0pt, fill=green!12, minimum height=0.82cm, inner xsep=7pt},
    link/.style={->, thick, draw=black!55},
    leaflink/.style={<->, thick, draw=red!65!black},
    promote/.style={->, dashed, very thick, draw=red!70!black}
]
    \node[font=\bfseries] at (7.6,8.2) {Before inserting 21};
    \node[inode] (r0) at (7.6,7.0) {5 \quad 13 \quad 17 \quad 24};

    \node[leaf] (a0) at (0.4,4.7) {2 \quad 3};
    \node[leaf] (b0) at (3.2,4.7) {5 \quad 7 \quad 8};
    \node[leaf] (c0) at (6.1,4.7) {14 \quad 16};
    \node[fullleaf] (d0) at (10.1,4.7) {19 \quad 20 \quad 22 \quad 23};
    \node[leaf] (e0) at (14.2,4.7) {24 \quad 27 \quad 29};

    \draw[link] (r0.south west) -- (a0.north);
    \draw[link] ($(r0.south west)!0.25!(r0.south east)$) -- (b0.north);
    \draw[link] ($(r0.south west)!0.50!(r0.south east)$) -- (c0.north);
    \draw[link] ($(r0.south west)!0.75!(r0.south east)$) -- (d0.north);
    \draw[link] (r0.south east) -- (e0.north);

    \draw[leaflink] (a0.east) -- (b0.west);
    \draw[leaflink] (b0.east) -- (c0.west);
    \draw[leaflink] (c0.east) -- (d0.west);
    \draw[leaflink] (d0.east) -- (e0.west);

    \node[font=\scriptsize, text=orange!80!black] at (10.1,3.9) {target leaf is full};

    \node[font=\bfseries] at (7.6,2.2) {After split propagation};
    \node[inode] (r1) at (7.6,1.0) {17};
    \node[inode] (i1) at (4.1,-1.0) {5 \quad 13};
    \node[inode] (i2) at (11.0,-1.0) {21 \quad 24};

    \node[leaf] (a1) at (0.3,-3.4) {2 \quad 3};
    \node[leaf] (b1) at (3.0,-3.4) {5 \quad 7 \quad 8};
    \node[leaf] (c1) at (5.9,-3.4) {14 \quad 16};
    \node[leaf] (d1) at (9.0,-3.4) {19 \quad 20};
    \node[newleaf] (e1) at (12.0,-3.4) {21 \quad 22 \quad 23};
    \node[leaf] (f1) at (15.1,-3.4) {24 \quad 27 \quad 29};

    \draw[link] (r1.south west) -- (i1.north);
    \draw[link] (r1.south east) -- (i2.north);

    \draw[link] (i1.south west) -- (a1.north);
    \draw[link] (i1.south) -- (b1.north);
    \draw[link] (i1.south east) -- (c1.north);

    \draw[link] (i2.south west) -- (d1.north);
    \draw[link] (i2.south) -- (e1.north);
    \draw[link] (i2.south east) -- (f1.north);

    \draw[leaflink] (a1.east) -- (b1.west);
    \draw[leaflink] (b1.east) -- (c1.west);
    \draw[leaflink] (c1.east) -- (d1.west);
    \draw[leaflink] (d1.east) -- (e1.west);
    \draw[leaflink] (e1.east) -- (f1.west);

    \draw[promote] (e1.north) to[bend left=15] node[midway, right, font=\scriptsize] {copy up 21} (i2.south);
    \draw[promote] (i1.north) to[bend left=18] node[midway, left, font=\scriptsize] {push up 17} (r1.south);
\end{tikzpicture}
}}\par}
\end{schema}

The sequence is:
\begin{dashlist}
    \item insert \(21\) into the full leaf \([19,20,22,23]\),
    \item split that leaf into \([19,20]\) and \([21,22,23]\),
    \item copy separator key \(21\) into the parent,
    \item the parent now overflows, so the parent must also split,
    \item push separator key \(17\) up to a brand-new root.
\end{dashlist}

The final tree is taller by one level. This is the standard cascading-insertion case: a local split at the leaf causes another split above it.

In theory, redistributing entries differently could delay a root split. In practice, most implementations use the simple split rule because it is local, predictable, and easy to maintain.

\subsection{Leaf split versus internal split}
One detail is especially important: \textbf{leaf splits} and \textbf{internal splits} do \emph{not} propagate keys in the same way.

\begin{dbtable}{Leaf split and internal split}{L{2.7cm}L{2.7cm}L{4.7cm}}
\textbf{Case} & \textbf{What moves upward} & \textbf{Reason} \\ \hline
Leaf split & separator key is \textbf{copied up} & The leaf stores actual data entries, so the key must remain in the leaf. The parent only needs a separator. \\ \hline
Internal split & middle separator is \textbf{pushed up} & Internal-node keys are only routing information, so the promoted separator appears only once, in the parent. \\
\end{dbtable}

This is the clean rule to remember:
\begin{dashlist}
    \item \textbf{leaf split} \(\Rightarrow\) \textbf{copy up},
    \item \textbf{internal split} \(\Rightarrow\) \textbf{push up}.
\end{dashlist}

\subsection{Another insertion sequence: 28, 6, and 25}
The next example shows three insertions in a row. Some fit locally. Others cause a split and change the upper levels as well.

\begin{schema}
{\resizebox{0.98\textwidth}{!}{\begin{tikzpicture}[
    x=0.95cm,y=0.87cm,
    every node/.style={font=\small},
    >=Latex,
    inode/.style={draw=black!60, rounded corners=3pt, line width=0.85pt, fill=blue!10, minimum height=0.82cm, inner xsep=9pt},
    leaf/.style={draw=black!60, rounded corners=3pt, line width=0.85pt, fill=blue!7, minimum height=0.82cm, inner xsep=7pt},
    link/.style={->, thick, draw=black!55},
    leaflink/.style={<->, thick, draw=red!65!black}
]
    \node[inode] (root) at (7.6,5.7) {13 \quad 30};
    \node[inode] (lefti) at (3.7,3.5) {5 \quad 7};
    \node[inode] (righti) at (11.7,3.5) {20 \quad 23};

    \node[leaf] (l1) at (0.5,1.1) {2 \quad 3};
    \node[leaf] (l2) at (3.1,1.1) {5 \quad 6};
    \node[leaf] (l3) at (5.7,1.1) {7 \quad 8 \quad 11};
    \node[leaf] (l4) at (9.1,1.1) {14 \quad 16};
    \node[leaf] (l5) at (11.8,1.1) {21 \quad 22};
    \node[leaf] (l6) at (15.0,1.1) {23 \quad 25 \quad 28};

    \draw[link] (root.south west) -- (lefti.north);
    \draw[link] (root.south east) -- (righti.north);

    \draw[link] (lefti.south west) -- (l1.north);
    \draw[link] (lefti.south) -- (l2.north);
    \draw[link] (lefti.south east) -- (l3.north);

    \draw[link] (righti.south west) -- (l4.north);
    \draw[link] (righti.south) -- (l5.north);
    \draw[link] (righti.south east) -- (l6.north);

    \draw[leaflink] (l1.east) -- (l2.west);
    \draw[leaflink] (l2.east) -- (l3.west);
    \draw[leaflink] (l3.east) -- (l4.west);
    \draw[leaflink] (l4.east) -- (l5.west);
    \draw[leaflink] (l5.east) -- (l6.west);

    \node[font=\scriptsize] at (7.6,-0.1) {final tree after inserting 28, 6, and 25};
\end{tikzpicture}
}}\par}
\end{schema}

Read the sequence like this:
\begin{dashlist}
    \item \textbf{insert 28}: the key fits into the rightmost leaf, so no split is needed,
    \item \textbf{insert 6}: the leaf \([5,7,8,11]\) overflows, so it splits into \([5,6]\) and \([7,8,11]\); key \(7\) is copied up,
    \item \textbf{insert 25}: the leaf \([21,22,23,28]\) overflows, so it splits into \([21,22]\) and \([23,25,28]\); key \(23\) is copied up,
    \item the parent then overflows, so an internal split pushes key \(13\) up and creates the final two-level tree shown above.
\end{dashlist}

This example reinforces the main idea: insertions start locally at one leaf, but the structural effect may travel upward.


\section{Deletion in B\texorpdfstring{$^+$}{+}-trees}
\begin{concept}{Deletion in a B\texorpdfstring{$^+$}{+}-tree}
Deletion removes a key from its leaf. If the leaf still satisfies the minimum-occupancy rule, the job is done. Otherwise the DBMS rebalances the tree by redistribution or merge.
\end{concept}

The standard deletion pattern is:
\begin{dashlist}
    \item start at the root and find the leaf \(L\) containing the entry,
    \item remove the entry from \(L\),
    \item if \(L\) is still at least half full, stop,
    \item otherwise try \textbf{redistribution}: borrow an entry from a sibling with the same parent,
    \item if no sibling can spare an entry, \textbf{merge} \(L\) with one sibling,
    \item after a merge, delete the corresponding separator entry from the parent,
    \item if the parent underflows, apply the same logic again higher up,
    \item if the root ends with only one child, remove that root and decrease the tree height.
\end{dashlist}

\subsection{Worked example: delete 19 and 20}
Deleting \(19\) is easy: after removing it, the leaf still has enough entries. Deleting \(20\) is more interesting, because the leaf becomes too small.

\begin{schema}
{\resizebox{0.98\textwidth}{!}{\begin{tikzpicture}[
    x=0.92cm,y=0.84cm,
    every node/.style={font=\small},
    >=Latex,
    inode/.style={draw=black!60, rounded corners=3pt, line width=0.85pt, fill=blue!10, minimum height=0.82cm, inner xsep=8pt},
    leaf/.style={draw=black!60, rounded corners=3pt, line width=0.85pt, fill=blue!7, minimum height=0.82cm, inner xsep=7pt},
    link/.style={->, thick, draw=black!55},
    leaflink/.style={<->, thick, draw=red!65!black}
]
    \node[font=\bfseries] at (7.4,8.0) {After deleting 19};
    \node[inode] (r0) at (7.4,6.8) {17};
    \node[inode] (i0l) at (3.7,4.8) {5 \quad 13};
    \node[inode] (i0r) at (11.2,4.8) {24 \quad 30};

    \node[leaf] (a0) at (0.5,2.4) {2 \quad 3};
    \node[leaf] (b0) at (3.1,2.4) {5 \quad 7 \quad 8};
    \node[leaf] (c0) at (5.9,2.4) {14 \quad 16};
    \node[leaf] (d0) at (10.0,2.4) {20 \quad 22};
    \node[leaf] (e0) at (12.8,2.4) {24 \quad 27 \quad 29};
    \node[leaf] (f0) at (16.2,2.4) {33 \quad 34 \quad 38 \quad 39};

    \draw[link] (r0.south west) -- (i0l.north);
    \draw[link] (r0.south east) -- (i0r.north);
    \draw[link] (i0l.south west) -- (a0.north);
    \draw[link] (i0l.south) -- (b0.north);
    \draw[link] (i0l.south east) -- (c0.north);
    \draw[link] (i0r.south west) -- (d0.north);
    \draw[link] (i0r.south) -- (e0.north);
    \draw[link] (i0r.south east) -- (f0.north);

    \draw[leaflink] (a0.east) -- (b0.west);
    \draw[leaflink] (b0.east) -- (c0.west);
    \draw[leaflink] (c0.east) -- (d0.west);
    \draw[leaflink] (d0.east) -- (e0.west);
    \draw[leaflink] (e0.east) -- (f0.west);

    \node[font=\bfseries] at (7.4,-0.2) {After deleting 20 and redistributing};
    \node[inode] (r1) at (7.4,-1.4) {17};
    \node[inode] (i1l) at (3.7,-3.4) {5 \quad 13};
    \node[inode] (i1r) at (11.2,-3.4) {27 \quad 30};

    \node[leaf] (a1) at (0.5,-5.8) {2 \quad 3};
    \node[leaf] (b1) at (3.1,-5.8) {5 \quad 7 \quad 8};
    \node[leaf] (c1) at (5.9,-5.8) {14 \quad 16};
    \node[leaf] (d1) at (10.0,-5.8) {22 \quad 24};
    \node[leaf] (e1) at (12.8,-5.8) {27 \quad 29};
    \node[leaf] (f1) at (16.2,-5.8) {33 \quad 34 \quad 38 \quad 39};

    \draw[link] (r1.south west) -- (i1l.north);
    \draw[link] (r1.south east) -- (i1r.north);
    \draw[link] (i1l.south west) -- (a1.north);
    \draw[link] (i1l.south) -- (b1.north);
    \draw[link] (i1l.south east) -- (c1.north);
    \draw[link] (i1r.south west) -- (d1.north);
    \draw[link] (i1r.south) -- (e1.north);
    \draw[link] (i1r.south east) -- (f1.north);

    \draw[leaflink] (a1.east) -- (b1.west);
    \draw[leaflink] (b1.east) -- (c1.west);
    \draw[leaflink] (c1.east) -- (d1.west);
    \draw[leaflink] (d1.east) -- (e1.west);
    \draw[leaflink] (e1.east) -- (f1.west);
\end{tikzpicture}
}}\par}
\end{schema}

After deleting \(19\), the leaf becomes \([20,22]\), which is still legal. After deleting \(20\), the same leaf becomes \([22]\), which violates minimum occupancy.

The DBMS therefore redistributes with the right sibling:
\begin{dashlist}
    \item the right sibling \([24,27,29]\) can spare one entry,
    \item the boundary moves, so the two leaves become \([22,24]\) and \([27,29]\),
    \item the parent separator is updated from \(24\) to \(27\).
\end{dashlist}

In a leaf-level redistribution, the new boundary key is \textbf{copied up} to the parent, because leaves keep the actual data entries.

\subsection{Worked example: delete 24 and shrink the tree}
Now continue from the tree above and delete \(24\). This time redistribution is no longer enough, so the tree must merge and then shrink.

\begin{schema}
{\resizebox{0.98\textwidth}{!}{\begin{tikzpicture}[
    x=0.92cm,y=0.84cm,
    every node/.style={font=\small},
    >=Latex,
    inode/.style={draw=black!60, rounded corners=3pt, line width=0.85pt, fill=blue!10, minimum height=0.82cm, inner xsep=8pt},
    leaf/.style={draw=black!60, rounded corners=3pt, line width=0.85pt, fill=blue!7, minimum height=0.82cm, inner xsep=7pt},
    link/.style={->, thick, draw=black!55},
    leaflink/.style={<->, thick, draw=red!65!black}
]
    \node[font=\bfseries] at (7.4,8.0) {Before deleting 24};
    \node[inode] (r0) at (7.4,6.8) {17};
    \node[inode] (i0l) at (3.7,4.8) {5 \quad 13};
    \node[inode] (i0r) at (11.2,4.8) {27 \quad 30};

    \node[leaf] (a0) at (0.5,2.4) {2 \quad 3};
    \node[leaf] (b0) at (3.1,2.4) {5 \quad 7 \quad 8};
    \node[leaf] (c0) at (5.9,2.4) {14 \quad 16};
    \node[leaf] (d0) at (10.0,2.4) {22 \quad 24};
    \node[leaf] (e0) at (12.8,2.4) {27 \quad 29};
    \node[leaf] (f0) at (16.2,2.4) {33 \quad 34 \quad 38 \quad 39};

    \draw[link] (r0.south west) -- (i0l.north);
    \draw[link] (r0.south east) -- (i0r.north);
    \draw[link] (i0l.south west) -- (a0.north);
    \draw[link] (i0l.south) -- (b0.north);
    \draw[link] (i0l.south east) -- (c0.north);
    \draw[link] (i0r.south west) -- (d0.north);
    \draw[link] (i0r.south) -- (e0.north);
    \draw[link] (i0r.south east) -- (f0.north);

    \draw[leaflink] (a0.east) -- (b0.west);
    \draw[leaflink] (b0.east) -- (c0.west);
    \draw[leaflink] (c0.east) -- (d0.west);
    \draw[leaflink] (d0.east) -- (e0.west);
    \draw[leaflink] (e0.east) -- (f0.west);

    \node[font=\bfseries] at (7.4,-0.2) {After merge and tree shrink};
    \node[inode] (r1) at (7.4,-1.5) {5 \quad 13 \quad 17 \quad 30};

    \node[leaf] (a1) at (0.5,-4.0) {2 \quad 3};
    \node[leaf] (b1) at (3.1,-4.0) {5 \quad 7 \quad 8};
    \node[leaf] (c1) at (5.9,-4.0) {14 \quad 16};
    \node[leaf] (d1) at (10.0,-4.0) {22 \quad 27 \quad 29};
    \node[leaf] (e1) at (14.2,-4.0) {33 \quad 34 \quad 38 \quad 39};

    \draw[link] (r1.south west) -- (a1.north);
    \draw[link] ($(r1.south west)!0.25!(r1.south east)$) -- (b1.north);
    \draw[link] ($(r1.south west)!0.50!(r1.south east)$) -- (c1.north);
    \draw[link] ($(r1.south west)!0.75!(r1.south east)$) -- (d1.north);
    \draw[link] (r1.south east) -- (e1.north);

    \draw[leaflink] (a1.east) -- (b1.west);
    \draw[leaflink] (b1.east) -- (c1.west);
    \draw[leaflink] (c1.east) -- (d1.west);
    \draw[leaflink] (d1.east) -- (e1.west);
\end{tikzpicture}
}}\par}
\end{schema}

The logic is:
\begin{dashlist}
    \item deleting \(24\) turns leaf \([22,24]\) into \([22]\), which is too small,
    \item its sibling \([27,29]\) is already at minimum occupancy, so borrowing is impossible,
    \item the two leaves merge into one leaf \([22,27,29]\),
    \item the parent loses one separator and now underflows,
    \item the merge therefore propagates upward,
    \item the old root is removed, and the tree height decreases by one.
\end{dashlist}

So deletion can do the opposite of insertion:
\begin{dashlist}
    \item insertion may \textbf{increase} the tree height,
    \item deletion may \textbf{decrease} the tree height.
\end{dashlist}

\subsection{Redistribution at non-leaf levels}
Redistribution is not only for leaves. It also works for internal nodes.

\begin{schema}
{\resizebox{0.98\textwidth}{!}{\begin{tikzpicture}[
    x=0.78cm,y=0.72cm,
    every node/.style={font=\scriptsize},
    >=Latex,
    inode/.style={draw=black!60, rounded corners=3pt, line width=0.8pt, fill=blue!10, minimum height=0.70cm, inner xsep=6pt},
    leaf/.style={draw=black!60, rounded corners=3pt, line width=0.8pt, fill=blue!7, minimum height=0.70cm, inner xsep=5pt},
    link/.style={->, thick, draw=black!55},
    leaflink/.style={<->, thick, draw=red!65!black}
]
    \node[font=\bfseries] at (8.5,8.5) {Before non-leaf redistribution};
    \node[inode] (r0) at (8.5,7.4) {22};
    \node[inode] (i0l) at (5.1,5.4) {5 \quad 13 \quad 17 \quad 20};
    \node[inode] (i0r) at (12.0,5.4) {27 \quad 30};

    \node[leaf] (a0) at (0.8,2.8) {2 \quad 3};
    \node[leaf] (b0) at (3.5,2.8) {5 \quad 7 \quad 8};
    \node[leaf] (c0) at (6.3,2.8) {14 \quad 16};
    \node[leaf] (d0) at (9.1,2.8) {17 \quad 18};
    \node[leaf] (e0) at (11.9,2.8) {20 \quad 21};
    \node[leaf] (f0) at (14.8,2.8) {22 \quad 24};
    \node[leaf] (g0) at (17.6,2.8) {27 \quad 29};
    \node[leaf] (h0) at (20.9,2.8) {33 \quad 34 \quad 38 \quad 39};

    \draw[link] (r0.south west) -- (i0l.north);
    \draw[link] (r0.south east) -- (i0r.north);

    \draw[link] (i0l.south west) -- (a0.north);
    \draw[link] ($(i0l.south west)!0.25!(i0l.south east)$) -- (b0.north);
    \draw[link] ($(i0l.south west)!0.50!(i0l.south east)$) -- (c0.north);
    \draw[link] ($(i0l.south west)!0.75!(i0l.south east)$) -- (d0.north);
    \draw[link] (i0l.south east) -- (e0.north);

    \draw[link] (i0r.south west) -- (f0.north);
    \draw[link] (i0r.south) -- (g0.north);
    \draw[link] (i0r.south east) -- (h0.north);

    \draw[leaflink] (a0.east) -- (b0.west);
    \draw[leaflink] (b0.east) -- (c0.west);
    \draw[leaflink] (c0.east) -- (d0.west);
    \draw[leaflink] (d0.east) -- (e0.west);
    \draw[leaflink] (e0.east) -- (f0.west);
    \draw[leaflink] (f0.east) -- (g0.west);
    \draw[leaflink] (g0.east) -- (h0.west);

    \node[font=\bfseries] at (8.5,-0.1) {After pushing entries through the parent};
    \node[inode] (r1) at (8.5,-1.2) {17};
    \node[inode] (i1l) at (5.1,-3.2) {5 \quad 13};
    \node[inode] (i1r) at (12.3,-3.2) {20 \quad 22 \quad 30};

    \node[leaf] (a1) at (0.8,-5.8) {2 \quad 3};
    \node[leaf] (b1) at (3.5,-5.8) {5 \quad 7 \quad 8};
    \node[leaf] (c1) at (6.3,-5.8) {14 \quad 16};
    \node[leaf] (d1) at (9.3,-5.8) {17 \quad 18};
    \node[leaf] (e1) at (12.1,-5.8) {20 \quad 21};
    \node[leaf] (f1) at (15.2,-5.8) {22 \quad 27 \quad 29};
    \node[leaf] (g1) at (19.0,-5.8) {33 \quad 34 \quad 38 \quad 39};

    \draw[link] (r1.south west) -- (i1l.north);
    \draw[link] (r1.south east) -- (i1r.north);

    \draw[link] (i1l.south west) -- (a1.north);
    \draw[link] (i1l.south) -- (b1.north);
    \draw[link] (i1l.south east) -- (c1.north);

    \draw[link] (i1r.south west) -- (d1.north);
    \draw[link] ($(i1r.south west)!0.33!(i1r.south east)$) -- (e1.north);
    \draw[link] ($(i1r.south west)!0.66!(i1r.south east)$) -- (f1.north);
    \draw[link] (i1r.south east) -- (g1.north);

    \draw[leaflink] (a1.east) -- (b1.west);
    \draw[leaflink] (b1.east) -- (c1.west);
    \draw[leaflink] (c1.east) -- (d1.west);
    \draw[leaflink] (d1.east) -- (e1.west);
    \draw[leaflink] (e1.east) -- (f1.west);
    \draw[leaflink] (f1.east) -- (g1.west);
\end{tikzpicture}
}}\par}
\end{schema}

The key idea is subtle but important:
\begin{dashlist}
    \item redistribution at the non-leaf level happens \textbf{through the parent},
    \item one separator from the parent moves down into the underfull internal node,
    \item at the same time, one separator from the fuller sibling moves up into the parent,
    \item the corresponding child pointer moves with that separator.
\end{dashlist}

So internal-node redistribution is not just ``move one key sideways''. It is really a \textbf{push-through-parent} operation.

\newpage
\section{Clustered indexes}
\begin{concept}{Clustered index}
An index whose logical order matches the physical order of the table data. Nearby index entries point to nearby records on disk.
\end{concept}

The main idea is physical:
\begin{dashlist}
    \item the table is stored in the sort order of one chosen key, often the primary key,
    \item the storage may be \textbf{heap-organized but kept in that order}, or \textbf{index-organized}, where the leaf level stores the actual rows,
    \item because the data file can have only one physical order at a time, a table can have at most one clustered organization.
\end{dashlist}

This is why clustered indexes are powerful for range scans:
\begin{dashlist}
    \item matching rows tend to sit on nearby pages,
    \item ordered scans can continue page by page with little random I/O.
\end{dashlist}

But there is also a cost:
\begin{dashlist}
    \item inserts and deletes are harder, because maintaining physical order may force page splits, movement, or reorganization.
\end{dashlist}

DBMSs differ here:
\begin{dashlist}
    \item some systems always organize a table around a clustered primary key; if no such key is declared, they create a hidden row identifier,
    \item other systems keep ordinary heap files and do not make the table itself follow one index order.
\end{dashlist}

\subsection{B\texorpdfstring{$^+$}{+}-tree traversal and scan order}
To scan many records in key order, the DBMS first descends to the leftmost relevant leaf and then follows leaf links from left to right.

\begin{schema}
{\begin{tikzpicture}[
    x=1cm,y=0.92cm,
    every node/.style={font=\small},
    >=Latex,
    frame/.style={draw=black!60, rounded corners=3pt, line width=0.85pt},
    leaf/.style={frame, fill=gray!10, minimum width=0.9cm, minimum height=0.55cm},
    page/.style={frame, fill=blue!7, minimum width=1.15cm, minimum height=0.72cm},
    flow/.style={->, thick, draw=red!70!black},
    toplink/.style={<->, thick, draw=red!65!black}
]
    \node[font=\bfseries] at (3.8,7.2) {Clustered scan};
    \node[font=\bfseries] at (3.8,0.3) {Unclustered scan};

    \node[leaf] (c1) at (1.0,5.8) {10};
    \node[leaf] (c2) at (2.3,5.8) {20};
    \node[leaf] (c3) at (3.6,5.8) {30};
    \node[leaf] (c4) at (4.9,5.8) {40};
    \node[leaf] (c5) at (6.2,5.8) {50};
    \draw[toplink] (c1.east) -- (c2.west);
    \draw[toplink] (c2.east) -- (c3.west);
    \draw[toplink] (c3.east) -- (c4.west);
    \draw[toplink] (c4.east) -- (c5.west);
    \draw[flow, very thick] (0.3,6.7) -- node[above, font=\scriptsize] {scan direction} (6.9,6.7);

    \node[page] (cp1) at (1.1,3.6) {101};
    \node[page] (cp2) at (2.8,3.6) {102};
    \node[page] (cp3) at (4.5,3.6) {103};
    \node[page] (cp4) at (6.2,3.6) {104};
    \draw[flow] (c1.south) to[out=-90,in=90] (cp1.north);
    \draw[flow] (c2.south) to[out=-90,in=90] (cp1.north);
    \draw[flow] (c3.south) to[out=-90,in=90] (cp2.north);
    \draw[flow] (c4.south) to[out=-90,in=90] (cp3.north);
    \draw[flow] (c5.south) to[out=-90,in=90] (cp4.north);

    \node[font=\scriptsize] at (3.7,2.9) {leaf order and table-page order are close};

    \node[leaf] (u1) at (1.0,1.6) {10};
    \node[leaf] (u2) at (2.3,1.6) {20};
    \node[leaf] (u3) at (3.6,1.6) {30};
    \node[leaf] (u4) at (4.9,1.6) {40};
    \node[leaf] (u5) at (6.2,1.6) {50};
    \draw[toplink] (u1.east) -- (u2.west);
    \draw[toplink] (u2.east) -- (u3.west);
    \draw[toplink] (u3.east) -- (u4.west);
    \draw[toplink] (u4.east) -- (u5.west);
    \draw[flow, very thick] (0.3,2.2) -- node[above, font=\scriptsize] {scan direction} (6.9,2.2);

    \node[page] (up1) at (1.1,-0.6) {101};
    \node[page] (up2) at (2.8,-0.6) {102};
    \node[page] (up3) at (4.5,-0.6) {103};
    \node[page] (up4) at (6.2,-0.6) {104};
    \draw[flow] (u1.south) to[out=-90,in=90] (up3.north);
    \draw[flow] (u2.south) to[out=-90,in=90] (up1.north);
    \draw[flow] (u3.south) to[out=-90,in=90] (up4.north);
    \draw[flow] (u4.south) to[out=-90,in=90] (up2.north);
    \draw[flow] (u5.south) to[out=-90,in=90] (up1.north);

    \node[font=\scriptsize] at (3.7,-1.5) {leaf order points to scattered table pages};
\end{tikzpicture}
}\par}
\end{schema}

The traversal itself is the same in both cases. The big difference is how the leaf entries map to table pages:
\begin{dashlist}
    \item with a \textbf{clustered} index, consecutive leaf entries usually point to nearby table pages, so scanning leaf after leaf also reads data pages in a mostly sequential order,
    \item with an \textbf{unclustered} index, consecutive leaf entries may point to tuples scattered over the table, so following the leaf order can trigger many repeated or random page reads.
\end{dashlist}

This explains why clustered tree indexes are excellent for ordered scans and large range queries.

For unclustered indexes, a common optimization is:
\begin{dashlist}
    \item first collect the matching RIDs or page IDs from the leaves,
    \item then sort or group them by \textbf{table page ID},
    \item finally fetch table pages in page order.
\end{dashlist}

This loses the natural key order during fetching, but it can save a lot of I/O when many tuples qualify.

\section{B\texorpdfstring{$^+$}{+}-tree design choices}
Real systems all use the same core \mbox{$B^+$-tree} idea, but they still make many low-level design choices.

\subsection{Node size}
Node size is one of the most important choices because it controls both:
\begin{dashlist}
    \item \textbf{fanout}: larger nodes make the tree shallower,
    \item \textbf{intra-node cost}: larger nodes take longer to search once loaded.
\end{dashlist}

As a rule of thumb, the slower the storage device, the larger the best node size tends to be:

\begin{dbtable}{Typical node-size tendencies}{L{2.8cm}C{2.2cm}L{4.2cm}}
\textbf{Storage medium} & \textbf{Typical size} & \textbf{Why} \\ \hline
HDD & \(\sim 1\) MB & Large transfers amortize expensive seeks and rotational delay. \\ \hline
SSD & \(\sim 10\) KB & Random access is cheaper, so extremely large pages are less necessary. \\ \hline
In-memory & \(\sim 512\) B & Cache behavior matters more than disk I/O, so smaller nodes can be better. \\
\end{dbtable}

These numbers are only rough guidelines. The best size still depends on the workload:
\begin{dashlist}
    \item long leaf scans benefit from larger nodes,
    \item repeated root-to-leaf lookups may prefer smaller nodes that fit caches better.
\end{dashlist}

\subsection{Merge threshold}
The textbook rule says that nodes should stay at least half full. In practice, many DBMSs do \emph{not} merge nodes immediately every time occupancy drops below that threshold.

Why delay merging?
\begin{dashlist}
    \item inserts are often more common than deletes, so a sparse page may soon fill again,
    \item eager merging causes extra page movement and separator updates,
    \item in practice, the average occupancy is often around \(67\%\),
    \item some systems prefer periodic rebuilding over aggressive merge-on-delete.
\end{dashlist}

So the theoretical half-full rule is best understood as a \textbf{guarantee for the ideal structure}, not always as a strict implementation policy after every update.

\subsection{Variable-length keys}
Variable-length search keys create another design problem: how do we store them inside fixed-size pages efficiently?

Common approaches include:
\begin{dashlist}
    \item \textbf{pointer-based keys}: store a pointer to the real attribute value instead of storing the full key inline,
    \item \textbf{variable-length nodes}: let the packed contents of one node vary in size and manage free space carefully,
    \item \textbf{padding}: extend each key to the maximum length of its type, which simplifies code but wastes space,
    \item \textbf{key indirection}: keep a small array of offsets or pointers that map to the actual key/value area inside the node.
\end{dashlist}

The trade-off is the usual one:
\begin{dashlist}
    \item simpler layout usually wastes more space,
    \item compact layout saves space but makes memory management harder.
\end{dashlist}

\subsection{Intra-node search}
Once one node is in memory, the DBMS still has to find the correct entry \emph{inside} that node.

\begin{dbtable}{Searching inside one node}{L{2.5cm}L{3.4cm}L{4.3cm}}
\textbf{Method} & \textbf{Main idea} & \textbf{When it is attractive} \\ \hline
Linear search & scan keys from left to right & good for small nodes; simple; often fast with SIMD/vectorized comparisons \\ \hline
Binary search & repeatedly compare with the middle key & good when nodes are larger and comparisons are cheap \\ \hline
Interpolation search & estimate the position from the key distribution & useful when keys are close to uniformly distributed \\
\end{dbtable}

For interpolation search, one common estimate is:
\[
\text{pos} \approx \text{low} +
\frac{(x - \text{arr}[\text{low}])(\text{high} - \text{low})}
{\text{arr}[\text{high}] - \text{arr}[\text{low}]},
\]
where \(x\) is the search key.

The idea is to jump close to the expected location instead of always probing the middle. This works best when key values are spread fairly evenly.

\Needspace{14\baselineskip}
\section{Implementation View: Nodes and Lookup}
\subsection{What one node stores}
At implementation level, a \mbox{$B^+$-tree} node is kept as a \textbf{sorted array} of entries inside one page.

The exact meaning of the value part depends on the node type:
\begin{dashlist}
    \item in an \textbf{internal node}, the value is a child page reference,
    \item in a \textbf{leaf node}, the value is a record pointer, a RID list, or the actual record itself, depending on the chosen data-entry alternative.
\end{dashlist}

Leaf nodes also keep \textbf{prev} and \textbf{next} links to neighboring leaves. These are what make range scans efficient.

Even though the internal-node layout has one extra rightmost child pointer, the high-level implementation view is still simple:
\begin{dashlist}
    \item every node stores entries in \textbf{sorted key order},
    \item search inside a node finds the first separator that tells us where to go next.
\end{dashlist}

\subsection{Lookup operation}
The lookup algorithm follows the sorted structure of the tree:
\begin{enumerate}
    \item Start at the root node.
    \item If the current node is an internal node, search inside it for the first separator that determines the correct child.
    \item Follow that child pointer to the next level.
    \item Repeat until a leaf is reached.
    \item In the leaf, either find the exact entry or determine the exact insertion position where the key would belong.
\end{enumerate}

So a lookup can return either:
\begin{dashlist}
    \item the \textbf{matching entry}, if the key exists,
    \item or the \textbf{correct leaf position}, if the key does not exist but might be inserted later.
\end{dashlist}

This lower-bound behavior is exactly what also makes insertions and range scans natural.

\section{Concurrent Access to B\texorpdfstring{$^+$}{+}-trees}
Concurrency is harder in trees than in isolated heap pages because pages depend on one another structurally.

Simple page latching is not enough:
\begin{dashlist}
    \item it protects one page against concurrent modification,
    \item but a split or merge also changes parent-child relationships,
    \item so several pages may need to be synchronized together.
\end{dashlist}

\begin{concept}{Lock coupling (latch crabbing)}
During a top-down traversal, the thread latches a parent, then latches the child it wants to enter, and only then releases the parent. The traversal therefore keeps a short overlapping chain of latches.
\end{concept}

This classical method works well for search and for many simple updates:
\begin{dashlist}
    \item latch the root first,
    \item latch the next level,
    \item release the parent,
    \item continue downward in canonical root-to-leaf order.
\end{dashlist}

Because latches are acquired in a fixed top-down order, ordinary traversals avoid deadlocks.

\subsection{Why inserts are trickier}
Insertion is harder than lookup because a split may propagate upward:
\begin{dashlist}
    \item a leaf split inserts a separator into the parent,
    \item that parent may also split,
    \item the split can continue all the way to the root.
\end{dashlist}

This creates a problem for naive lock coupling: by the time the split reaches higher levels, those ancestors may already have been released.

\subsection{Restart or optimistic coupling}
A simple practical solution is:
\begin{enumerate}
    \item first try the insertion using ordinary lock coupling,
    \item if no split propagates through inner nodes, the operation succeeds normally,
    \item otherwise release the latches,
    \item restart the operation with a more conservative locking strategy that keeps the needed ancestors protected,
    \item execute the update safely on the second pass.
\end{enumerate}

This approach has a clear trade-off:
\begin{dashlist}
    \item it greatly reduces concurrency during the restarted operation,
    \item but the expensive case is rare,
    \item and the algorithm is much simpler to implement correctly.
\end{dashlist}

\section{B\texorpdfstring{$^+$}{+}-trees in Practice}
The main practical lesson is that \mbox{$B^+$-trees} are usually \textbf{very shallow}.

Many textbooks parameterize them with a value \(d\), where an internal node has between \(d\) and \(2d\) children. If we take a typical example with:
\begin{dashlist}
    \item \(d = 100\),
    \item average fill factor \(\approx 67\%\),
\end{dashlist}
then the average fanout is roughly
\[
2d \cdot 0.67 = 200 \cdot 0.67 \approx 134.
\]

That already gives very large capacities:

\begin{dbtable}{Back-of-the-envelope capacities}{L{2.8cm}L{4.8cm}}
\textbf{Tree height} & \textbf{Approximate number of entries} \\ \hline
3 & \(134^3 = 2{,}406{,}104\) \\ \hline
4 & \(134^4 = 322{,}417{,}936\) \\
\end{dbtable}

This is why a tree of height \(3\) or \(4\) is often enough even for very large relations.

Top levels are also small enough to remain in memory:

\begin{dbtable}{Example memory footprint with 8 KB pages}{L{2.5cm}L{4.8cm}}
\textbf{Level} & \textbf{Size} \\ \hline
level 1 & \(1\) page \(\approx 8\) KB \\ \hline
level 2 & \(134\) pages \(\approx 1\) MB \\ \hline
level 3 & \(17{,}956\) pages \(\approx 140\) MB \\
\end{dbtable}

So in practice:
\begin{dashlist}
    \item the root is almost always memory-resident,
    \item often one or more upper levels are memory-resident too,
    \item only the lower levels require regular I/O.
\end{dashlist}

\end{document}
